% !TeX root = ../prompt-engineering-guide-for-beginners.tex

%%%%%%%%%%%%%%%%%%%%%%%%%%%%%%%%%%
%% 第五章：设计提示的通用技巧
%%%%%%%%%%%%%%%%%%%%%%%%%%%%%%%%%%
\chapter{设计提示的通用技巧：写出高分 Prompt 的方法论}
\label{ch:general-tips}

上一章我们拆解了提示词的结构，现在来聊一些经过大量实践验证的实用技巧。
这些技巧不针对某个特定场景，而是适用于所有类型的提示词。
掌握了这些方法论，你写提示词的水平会有一个明显的跃升。

\section{动词要具体，拒绝假大空}

很多人写提示词的习惯是用「说一下」「帮我看看」「讲一讲」这种笼统的动词。AI 拿到这样的指令，往往不知道你想要什么样的输出，只能给一个泛泛的回答。

解决办法很简单：换成更精确的动词。

\begin{center}
\begin{tabular}{ll}
  \toprule
  \textbf{模糊动词} & \textbf{精确动词} \\
  \midrule
  说一下 & 列出、总结、解释、分析 \\
  帮我看看 & 检查、对比、评估、纠错 \\
  讲一讲 & 描述、定义、举例说明 \\
  弄一下 & 改写、翻译、格式化、提炼 \\
  \bottomrule
\end{tabular}
\end{center}

例如：

\begin{itemize}
  \item 「说说这个产品的优缺点」→「\textbf{列出}这个产品的 3 个优点和 3 个缺点，每个配一句简短的解释」
  \item 「帮我看看这段代码」→「\textbf{检查}这段代码是否存在性能问题，并给出优化建议」
\end{itemize}

一个更精确、聚焦的动词，就能让 AI 瞬间领会侧重点，从而将回答质量提升一整个档次。

\section{提供上下文：AI 不会读心术}

AI 只能根据你告诉它的信息来回答。你脑子里知道的背景，如果没有写进提示词，AI 就是不知道。

所以，写提示词的时候要养成一个习惯：把相关的背景信息尽量提供给 AI。

\begin{itemize}
  \item 你是谁（角色/身份）
  \item 你的受众是谁（给谁看/用的）
  \item 你的目标是什么（想达到什么效果）
  \item 有什么限制条件（公司规定、预算、时间等）
  \item 有没有参考资料（是的话，直接贴进来）
\end{itemize}

\textbf{反面例子：}
\begin{promptbox}
帮我写一段欢迎词。
\end{promptbox}

\textbf{正面例子：}
\begin{promptbox}
我是一家少儿编程培训机构的校长，下周六有一场面向 6 到 12 岁小朋友家长的开放日活动。请帮我写一段 3 分钟的欢迎词，语气热情但不夸张，重点介绍我们的教学特色和孩子们能收获什么。
\end{promptbox}

信息越是生动丰满，AI 的回答就越像是一个深入一线、懂你业务的资深同行执笔的，而不是一个走过场敷衍了事的机器路人。

\section{设定约束条件}

不给约束条件的提示词就像是一张没有边界的白纸，AI 不知道该往哪个方向写，只能写一大堆「正确但没用」的内容。

常见的约束条件包括：

\begin{itemize}
  \item \textbf{字数限制}：「控制在 200 字以内」「不超过 3 段」
  \item \textbf{语气要求}：「正式/轻松/幽默/严肃」
  \item \textbf{受众定义}：「写给完全没有技术背景的人看」
  \item \textbf{排除条件}：「不要使用专业术语」「不要列举超过 5 个」
  \item \textbf{格式限定}：「以 Markdown 表格形式输出」「用编号列表」
\end{itemize}

约束条件越是详细，AI 的产物就越贴合你的心意。你可以把约束条件想象成预先给 AI 画好了一个边界分明的工作框：在这个框以内，它拥有极大地文字组织自由；但框外那些「正确但你不需要」的废话，统统被强制过滤掉了。

\section{分步拆解复杂任务}

如果面对的需求本身异常庞大复杂，试图用“一言以蔽之”的方式把所有指令全塞进对话框里是非常危险的，往往会导致 AI 在生成时顾此失彼。
针对此类场景，最稳妥的策略就是化繁为简：把一个超级任务肢解为先后承接的几个小短步。

\textbf{不太好的做法：}
\begin{promptbox}
帮我写一篇关于远程办公趋势的深度文章，要有数据引用、案例分析、未来展望，3000 字以上。
\end{promptbox}

\textbf{更好的做法：}

\begin{promptbox}
第一步：请帮我列出关于「远程办公趋势」的文章大纲，包括 5 到 7 个主要章节。

第二步（拿到大纲并确认没问题后）：请基于以上大纲的指引，逐章推进撰写。请先集中精力只完成第一章的内容：远程办公的现状与底层数据支撑。

第三步（如此循环，直到全文章节完成）：现全文已生成完毕，请你作为最终主编通读全文，集中修正可能存在的逻辑断层，并专门润色章节之间的过渡段，使全文阅读如同行云流水。
\end{promptbox}

分步拆解最大的优势在于：每个独立小步骤的任务边界都极其明确清晰，AI 在局部的正确率和表现力上就会达到峰值。更重要的是，在每一轮微小动作交付后，你都能及时进行干预修正，而不必面临“等了 3 千字出来却发现必须全盘推翻重写”的惨痛沉没成本。

\section{使用分隔符组织内容}

当你的提示词里既有指令又有输入数据时，用分隔符把它们清晰地隔开，AI 就不会搞混。

常用的分隔符有：

\begin{itemize}
  \item 三个引号：\verb|"""|
  \item 三个反引号：\verb|```|
  \item XML 标签：\verb|<text>...</text>|
  \item 横线：\verb|---|
\end{itemize}

\textbf{不用分隔符的例子：}
\begin{promptbox}
请总结以下内容的核心观点远程办公正逐渐成为许多企业的常态化选择根据最新调查显示超过60\%的员工希望每周至少有两天在家办公……
\end{promptbox}

AI 可能分不清哪些是你的指令、哪些是需要处理的内容。

\textbf{用了分隔符的例子：}
\begin{promptbox}
请总结以下内容的核心观点：

\verb|"""|

远程办公正逐渐成为许多企业的常态化选择。根据最新调查显示，超过 60\% 的员工希望每周至少有两天在家办公……

\verb|"""|
\end{promptbox}

用分隔符把指令和数据区分开后，AI 就能清楚地知道哪些是你的操作指令、哪些是需要处理的原始素材，从而避免混淆。

\section{先说「做什么」而非「不做什么」}

人在提需求时，常常会习惯性地使用否定句式来做限制，比如「不要写得太长」「不要使用专业术语」「不要太正式」。

这里面有个有趣的现象：当你告诉 AI「不要写得太长」时，模型内部反而会先激活与「太长」相关的概念，这可能导致它偏向那些你想要避免的方向。相比之下，直接用正向指令（明确说出你想要什么）的效果通常更好。道理很简单：你告诉了 AI 不要做什么，但它仍然不知道你到底想要什么。

\begin{center}
\begin{tabular}{ll}
  \toprule
  \textbf{负向指令} & \textbf{正向指令} \\
  \midrule
  不要写得太长 & 控制在 150 字以内 \\
  不要用专业术语 & 用通俗易懂的语言 \\
  不要太正式 & 用轻松的口语化语气 \\
  不要列举太多 & 只列出最重要的 3 点 \\
  \bottomrule
\end{tabular}
\end{center}

当然，绝对性禁令在某些底线场合依旧是无可替代的（例如系统指令里规定的「绝对禁止编造虚假数据并以此对用户作答」）。但是在 90\% 的日常指令沟通下，试着切换心态，先把「鼓励它走向什么模样」描绘得清清楚楚，成效必然会截然不同。如果它仍然有跑偏的迹象，你再在句尾补上一刀轻微的反向封禁即可。

到这里，你已经掌握了写好提示词的通用方法论。接下来的三章，我们将按照「入门 $\to$ 进阶 $\to$ 高阶」的梯度，为你介绍提示工程领域中几种经过验证的具体技巧。别被名字吓到，其中很多方法你可能已经在实践中无意识地用过了。
